\section{总结与延伸}

本节把整期压缩成一条路线。第一，核能来自质量亏损，可控核聚变要把不可控释放变成稳定能源产品。第二，托卡马克是磁约束中最有积累的路线，高温超导材料把路线的关键从“能不能约束”推进到“能不能低成本做小”。第三，三重积和 Q 值解释了物理门槛，度电成本解释了商业门槛。第四，创业公司适合在科学可行性之后压缩工程和成本，但需要极强资本耐心。第五，AI 和能源会互相塑造：AI 需要更多电，聚变也需要 AI 降本增效。

\begin{importantbox}{把本期放进张小珺 AI/互联网队列}
这期虽然不是 AI 公司访谈，但它解释了 AI 时代的底层能源约束。与 EP110/EP115 的 Agent 系统工程、EP106/EP109 的物理世界数据路线相互补充：如果数字智能要继续放大，物理能源和真实世界工程会重新成为关键变量。
\end{importantbox}

\subsection{关键 takeaways}

\begin{enumerate}
\item 可控核聚变的难点不是“聚变能否发生”，而是能否稳定、可控、低成本地发电。
\item 高温超导托卡马克的核心价值是高磁场带来的小型化和成本下降。
\item Q 值和三重积是物理性能指标，度电成本才是商业化金标准。
\item 洪荒 70 验证完整高温超导装置，洪荒 170 指向 Q 大于 10，洪荒 380 指向示范电站产品。
\item AI 和核聚变互为条件：AI 增加能源需求，AI 技术也能用于聚变控制、诊断和仿真。
\end{enumerate}

\subsection{拓展阅读}

\begin{itemize}
\item 想理解 AI 与能源约束，可对照 EP127/EP136 中关于数据中心、GPU/TPU 和模型公司资本开支的讨论。
\item 想理解物理世界工程创业，可对照 EP111 激光雷达创业、EP106/EP109 机器人生产力和仿真数据路线。
\item 想把这期当科普课复习，优先掌握三个公式/概念：$E=\Delta mc^2$、三重积 $nT\tau_E$、能量增益 $Q=P_{\mathrm{fusion}}/P_{\mathrm{input}}$。
\end{itemize}

\end{document}
